\documentclass[11pt]{article}
\usepackage[T1]{fontenc}
\usepackage{lmodern,amsmath,amssymb,amsthm,mathtools}
\usepackage[margin=1in]{geometry}
\usepackage[hidelinks]{hyperref}
\newtheorem{theorem}{Theorem}[section]
\newtheorem{lemma}[theorem]{Lemma}
\newtheorem{proposition}[theorem]{Proposition}
\newcommand{\R}{\mathbb R}
\newcommand{\E}{\mathbb E}
\newcommand{\HH}{\mathcal H}
\newcommand{\DD}{\mathcal D}
\newcommand{\Smo}{\mathsf S}
\newcommand{\Reg}{\mathsf A}
\newcommand{\J}{\mathsf J}
\newcommand{\code}[1]{{\small\nolinkurl{#1}}}
\title{A common number regularization for\protect\\Brownian clock calculus:\protect\\full predictable limits and correction duality}
\author{Project Stardust\\\small A formalization companion}
\date{Reviewed editorial draft, October 30, 2026}
\begin{document}
\maketitle
\begin{abstract}
Let $W$ be a two-sided Brownian scenery independent of a Brownian clock $B$.
For a real Schwartz function $g$, spatially smooth versions of
$g(a+W_{|B_t|})-g(a)$ have a genuine scenery-valued reflected It\^o formula.
We prove that one fixed injective smoothing of the clock Gaussian number
operator gives a strong limit of the entire predictable integrand in the
original time--clock--scenery $L^2$ space. The full curvature and local-time
correction converges in the associated common dual norm. The limits are
independent of the fixed nonnegative compact $C^3$ mass-one profile, and
hold as its positive width tends continuously to zero. The decisive bound
is proportional to $t^{-1/2}$ and is integrable at the origin. Fixed clock
sectors converge in ordinary joint $L^2$, while independent rectangular
source/clock truncations have uniformly small tails in the dual norm.
Every statement retains the original Brownian-clock filtration and all
scenery variables. No unsmoothed correction limit or intrinsic
observed-filtration IBM It\^o formula is asserted. Classical Gaussian
spectral and Hilbert-valued stochastic mechanisms are credited; research
priority for this particular assembly remains unresolved.
\end{abstract}

\section{The original spaces and the result}
Fix independent standard Brownian sources $W$ and $B$, with
$W$ two-sided and $B_0=0$. Thus
\[
 \E_W(W_r-W_s)^2=|r-s|,\qquad \E_B B_t^2=t.
\]
The reflected-clock observation is $Z_t=W_{|B_t|}$. Work with real spaces
\[
 E=L^2(\Omega_W),\qquad
 \HH=L^2(\Omega_B;E)\simeq L^2(\Omega_B\times\Omega_W).
\]
All these are the original sources, with their full Gaussian sigma
algebras. Fubini identifications are isometries and do not replace the
joint variables by variables with matching laws. A scenery-valued
predictable process means a strongly measurable $E$-valued process
predictable for the Brownian clock filtration. Its space at horizon
$T\geq0$ is the closed Hilbert space
\[
 \DD_T=L^2_{\mathrm{pred}}((0,T]\times\Omega_B;E).
\]
The original Hilbert-valued It\^o integral is the isometry
\begin{equation}\label{eq:ito}
 \J_T:\DD_T\longrightarrow\HH,\qquad
 \J_Tu=\int_0^T u_t\,dB_t,\qquad \|\J_Tu\|_{\HH}=\|u\|_{\DD_T}.
\end{equation}
The scenery is the value-space parameter; it is not averaged out.
Equivalently one may reveal the whole independent scenery at time zero
and integrate in $B$. This convention is different from predictability
for the smaller observed past $\sigma(Z_s:s\leq t)$.

Let $P_k$ denote the orthogonal projection onto the $k$th clock Wiener
chaos, tensored with the identity on $E$. Put
\begin{equation}\label{eq:smooth}
 \Smo=P_0+N_B^{-1},\qquad
 N_B^{-1}f=\sum_{k\geq1}\frac1k P_kf.
\end{equation}
The inverse here is zero on the clock vacuum. The added $P_0$ is
essential. In particular $\Smo$ is an injective self-adjoint contraction
on $\HH$, retaining all scenery chaoses.

Fix $g\in\mathcal S(\R;\R)$ and $a\in\R$. A profile is a fixed function
$\rho\in C_c^3(\R)$ such that $\rho\geq0$ and $\int\rho=1$.
Evenness is not required. For $\varepsilon>0$ set
\begin{align}
 \rho_\varepsilon(x)&=\varepsilon^{-1}\rho(x/\varepsilon),\nonumber\\
 S_\varepsilon(r)&=\int_\R\rho_\varepsilon(r-y)W_y\,dy\quad\hbox{in }E,\label{eq:scenery}\\
 h_\varepsilon(r)&=g(a+S_\varepsilon(r)),\nonumber\\
 F_{\varepsilon,t}&=h_\varepsilon(|B_t|)-h_\varepsilon(0),\qquad
 F_{0,t}=g(a+Z_t)-g(a).\label{eq:endpoint}
\end{align}
The integral in \eqref{eq:scenery} is a Bochner convolution. The same
$S_\varepsilon$ is used at both endpoint arguments. The random
$h_\varepsilon(0)$ generally differs from $g(a)$.

The raw integrand and raw correction are
\begin{align}
 U_{\varepsilon,t}
  &=\operatorname{sgn}(B_t)h_\varepsilon'(|B_t|),\label{eq:rawu}\\
 C_{\varepsilon,T}
  &=L_T^0(B)h_\varepsilon'(0)
       +\frac12\int_0^T h_\varepsilon''(|B_t|)\,dt.\label{eq:rawc}
\end{align}
Local time has the Tanaka convention
$|B_T|=\int_0^T\operatorname{sgn}(B_t)\,dB_t+L_T^0(B)$.
For each positive width these are actual elements of $\DD_T$ and
$\HH$, respectively. Define the bounded predictable operator
\begin{equation}\label{eq:A}
 \Reg_Tu=\sum_{k\geq0}\frac1{k+1}\Pi_ku,
\end{equation}
where $\Pi_k$ is clock-chaos projection on the genuine time sections.
Its construction, predictability and integral intertwining are proved below.

\begin{theorem}[Full common regularized limit]\label{thm:main}
For every $g,a,T$ as above there are $u_{0,T}\in\DD_T$ and
$c_{0,T}\in\HH$, independent of the fixed admissible profile, such that
as $\varepsilon\downarrow0$ through all positive real widths,
\begin{align}
 \Reg_T U_\varepsilon&\longrightarrow u_{0,T}
                    &&\text{in }\DD_T,\label{eq:uconv}\\
 \Smo C_{\varepsilon,T}&\longrightarrow c_{0,T}
                    &&\text{in }\HH,\label{eq:cconv}\\
 \Smo F_{0,T}&=\J_Tu_{0,T}+c_{0,T}.\label{eq:limitidentity}
\end{align}
For $t>0$, with $v_t=1_{(0,t]}/\sqrt t$ in the clock Cameron--Martin
space, define $R_v=\langle v,D_BN_B^{-1}\,\cdot\,\rangle$. The time
sections of the limit are
\begin{equation}\label{eq:qzero}
 u_{0,T}(t)=t^{-1/2}R_{v_t}F_{0,t}
 \quad\text{for almost every }0<t\leq T.
\end{equation}
For $L=\sup|g'|$ and $m_1=\E|G|$ for a standard real Gaussian,
\begin{equation}\label{eq:energy}
 \|u_{0,T}(t)\|_{\HH}^2\leq L^2m_1t^{-1/2},\qquad
 \|u_{0,T}\|_{\DD_T}^2\leq2L^2m_1\sqrt T.
\end{equation}
The same section bound holds for $\Reg_TU_\varepsilon$, uniformly in
the width and the admissible profile.
\end{theorem}
For $T=0$ all spaces, integrals and endpoint differences in this
assertion are zero. No statement about a common exceptional sample set
for all terminal times is needed: this is a theorem at each fixed horizon.

\paragraph{Verified scope.}
The frozen source is commit
\par\noindent \code{bab688bb6140cd6316c80680095a914ea0d346db}.
The accompanying theorem map records exact declaration headers, source
hashes and permanent line links. The principal endpoint, predictable
limit and dual-closure modules have been reviewed with their proofs.
The existing full-build and complete transitive axiom-audit receipts
are checked by hash; this editorial paper changes no Lean source and
does not claim a new library build. Gaussian chaos completeness,
the constructed original Brownian It\^o isometry, and its established
Tanaka/local-time formula are prerequisites. We reconstruct their
application and the convergence proof, rather than relabeling those
prerequisites as new theorems.

\section{Spatial smoothing without a hidden profile constant}
The map $r\mapsto W_r$ is continuous into $E$, with exact increment
norm $\|W_r-W_s\|_E=\sqrt{|r-s|}$. Hence the convolution in
\eqref{eq:scenery} exists on a compact interval for every $r$.
Changing variables gives the useful identity
\[
 S_\varepsilon(r)=\int_\R\rho(z)W_{r-\varepsilon z}\,dz.
\]
This change of variables also fixes the sign convention for nonsymmetric
profiles.

\begin{lemma}[Uniform increment and approximation bounds]\label{lem:spatial}
Suppose $\operatorname{supp}\rho\subset[-R,R]$, $R\geq0$. Then
\begin{align}
 \|S_\varepsilon(r)-S_\varepsilon(s)\|_E
      &\leq\sqrt{|r-s|},\label{eq:inc}\\
 \|S_\varepsilon(r)-W_r\|_E^2&\leq\varepsilon R,\label{eq:error}\\
 \|h_\varepsilon(r)-h_\varepsilon(s)\|_E
      &\leq L\sqrt{|r-s|}.\label{eq:hinc}
\end{align}
In particular $h_\varepsilon(r)\to g(a+W_r)$ in $E$ at every fixed $r$.
\end{lemma}
\begin{proof}
Minkowski's integral inequality, positivity of $\rho$ and its unit mass
give
\[
 \Big\|\int\rho(z)(W_{r-\varepsilon z}-W_{s-\varepsilon z})\,dz\Big\|_E
 \leq\int\rho(z)\sqrt{|r-s|}\,dz=\sqrt{|r-s|}.
\]
For the approximation subtract $W_r\int\rho(z)\,dz$ inside the
Bochner integral and use $\sqrt{\varepsilon|z|}\leq\sqrt{\varepsilon R}$.
The ordinary mean value theorem gives
$|g(a+x)-g(a+y)|\leq L|x-y|$. Apply it to the genuine jointly
Gaussian scenery variables and take their $L^2$ norm. This proves
\eqref{eq:hinc} and the last assertion.
\end{proof}
Neither derivative norms of $\rho_\varepsilon$ nor a bound on $R$
enters \eqref{eq:inc} or \eqref{eq:hinc}. Radius $R$ occurs only in
the approximation estimate. The proof uses a fixed admissible profile;
it does not assert convergence for arbitrary width-dependent profiles
with unbounded supports or unproved approximation properties.

\begin{lemma}[Endpoint convergence and energy]\label{lem:endpoint}
For every fixed $t\geq0$, $F_{\varepsilon,t}\to F_{0,t}$ in $\HH$.
For all positive widths,
\begin{equation}\label{eq:Fenergy}
 \|F_{\varepsilon,t}\|_{\HH}^2\leq L^2m_1\sqrt t,
 \qquad \|F_{0,t}\|_{\HH}^2\leq L^2m_1\sqrt t.
\end{equation}
\end{lemma}
\begin{proof}
For every deterministic $r$, Lemma~\ref{lem:spatial} gives
$h_\varepsilon(r)-h_\varepsilon(0)\to g(a+W_r)-g(a)$ in $E$.
Both differences have $E$ norm at most $2\sup|g|$. Integrating their
squared difference against the law of $|B_t|$ proves joint $L^2$
convergence by dominated convergence. Alternatively the uniform
scenery error gives a quantitative endpoint bound by two applications
of the Lipschitz inequality. For the energy, first hold the clock
fixed and use \eqref{eq:hinc} with $s=0$. Its squared norm is at most
$L^2|B_t|$. Since $\E|B_t|=m_1\sqrt t$, Fubini proves the first
bound. Take the norm limit for the second. At $t=0$ the endpoint
differences are identically zero.
\end{proof}

\par\noindent Source correspondences:
\par\noindent \code{IBMContinuousScenery.sourceScenery_increment_sq},
\par\noindent \code{IBMDilatedMollifier.scenery_error},
\par\noindent \code{IBMCompactKernelIncrements.spatial_increment},
\par\noindent \code{IBMCompactClockGraph.endpoint_uniform_energy},
and \par\noindent \code{IBMDilatedEndpointLimit.endpoint_continuum_tendsto}.

\section{The genuine finite-width reflected formula}
Finite-width regularity is used before taking any singular limit.
Differentiation of the compact convolution in $E$ gives $S_\varepsilon'$
and $S_\varepsilon''$. Gaussian moment estimates and bounded derivatives
of $g$ justify the $E$-valued chain rules
\begin{align}
 h_\varepsilon'(r)&=g'(a+S_\varepsilon(r))S_\varepsilon'(r),\label{eq:hprime}\\
 h_\varepsilon''(r)&=g''(a+S_\varepsilon(r))(S_\varepsilon'(r))^2
                  +g'(a+S_\varepsilon(r))S_\varepsilon''(r).\label{eq:hsecond}
\end{align}
Here the products are actual scenery variables. They are not defined
by assigning Gaussian coefficients after the fact. Indeed the
spatial derivative of smoothed Brownian scenery has variance
$\|\rho_\varepsilon\|_2^2$, and the second derivative has variance
$\|\rho_\varepsilon'\|_2^2$. Thus globally in $r$,
\begin{align*}
 \|h_\varepsilon'(r)\|_E
     &\leq L\|\rho_\varepsilon\|_2,\\
 \|h_\varepsilon''(r)\|_E
     &\leq 2\sup|g''|\,\|\rho_\varepsilon\|_2^2
                +L\|\rho_\varepsilon'\|_2.
\end{align*}
The factor $2$ bounds the square root of the Gaussian fourth moment
factor $3$. Compact $C^3$ regularity ensures the required continuity
of the second derivative. To justify the chain rule directly, use
Gaussian $L^p$ bounds for the convolution difference quotients for
each finite $p$, Taylor's formula with the bounded derivatives of
$g$, and H\"older's inequality for the remainder products. Their
limits in $L^2$ give \eqref{eq:hprime}--\eqref{eq:hsecond}.

\begin{proposition}[Original finite-width identity]\label{prop:finite}
For every $\varepsilon>0$, $U_\varepsilon\in\DD_T$,
$C_{\varepsilon,T}\in\HH$, and
\begin{equation}\label{eq:finite}
 F_{\varepsilon,T}=\J_TU_\varepsilon+C_{\varepsilon,T}.
\end{equation}
\end{proposition}
\begin{proof}
Continuity of the $E$-valued slope and Brownian predictability show
that \eqref{eq:rawu} has a predictable version. The preceding global
slope bound makes its time $L^2$ norm finite. The curvature term in
\eqref{eq:rawc} is a pathwise Bochner time integral of a continuous
$E$-valued function and has deterministically bounded $E$ norm for
this fixed width. The boundary term belongs to $\HH$ because
$h_\varepsilon'(0)\in E$ and $L_T^0(B)\in L^2(\Omega_B)$.

For any continuous real linear functional $\ell$ on $E$, apply the
scalar $C^2$ It\^o formula to $\ell(h_\varepsilon(|B_t|))$.
Tanaka gives $d|B_t|=\operatorname{sgn}(B_t)dB_t+dL_t^0$;
its quadratic variation is $dt$, since the zero set has Lebesgue
measure zero almost surely. Local-time mass is supported on
$B_t=0$, so the supported Stieltjes integral is
$\ell(h_\varepsilon'(0))L_T^0$. The remaining finite-variation
term is exactly half the curvature integral. All scalarizations
commute with the original Hilbert-valued It\^o and Bochner integrals.
This proves \eqref{eq:finite} after every $\ell$.
Choose a countable separating family in the separable Hilbert space
$E$ and intersect its full-probability sets to recover the
$E$-valued identity. It is an equality in $\HH$.
\end{proof}
In particular \eqref{eq:rawc} constructs the correction before
\eqref{eq:finite} is used. The expression is not merely a name for
an unspecified residual. Individual derivative bounds here grow
as the width shrinks, so they do not prove convergence of either
raw term.

\par\noindent Source correspondences:
\par\noindent \code{IBMCompactKernelSchwartz.spatial_hasDerivAt},
\par\noindent \code{IBMCompactKernelSchwartz.slope_hasDerivAt}, and the actual
integrand, whole correction and finite identity in
\par\noindent \code{IBMCompactKernelReflectedIto}. The constructed Tanaka
and supported-local-time theorem is the prerequisite also used
in the sticky-clock and one-time-law companions~\cite{companions}.

\section{Clock number smoothing and predictable degree shift}
We give the operator mechanisms explicitly. Orthogonality yields
$f=\sum_{k\geq0}P_kf$ strongly in $\HH$. With
\[
 w_0=1,\qquad w_k=1/k\quad(k\geq1),
\]
\eqref{eq:smooth} is the strongly convergent sum $\sum w_kP_k$.
Its norm is at most one. The real weights prove self-adjointness.
If $\Smo f=0$, projection gives $w_kP_kf=0$ for every $k$;
all weights are positive, so $f=0$. Its range is dense: finite
clock-chaos sums have preimages obtained by multiplying their
finitely many components by $w_k^{-1}$, and these sums are dense.
These arguments apply to arbitrary $E$-valued coefficients and
retain all scenery sectors.

\begin{lemma}[Predictable projection and shift]\label{lem:shift}
Clock-chaos projection on time sections preserves $\DD_T$ and defines
an orthogonal projection $\Pi_k$ there. For every $u\in\DD_T$,
\begin{equation}\label{eq:shift}
 \J_T\Pi_ku=P_{k+1}\J_Tu.
\end{equation}
The strong series \eqref{eq:A} defines a contraction on $\DD_T$ and
\begin{equation}\label{eq:intertwine}
 \J_T\Reg_Tu=N_B^{-1}\J_Tu=\Smo\J_Tu.
\end{equation}
\end{lemma}
\begin{proof}
Start with scalar predictable elementary fields
$1_{(s,t]}F$, where $F$ is a polynomial of finitely many Brownian
coordinates from times at most $s$. Its chaos projection is still
a function of those past coordinates, hence is predictable.
More generally conditional expectation onto the Brownian past
is second quantization of the orthogonal projection onto
$L^2(0,s]$; applying it to each homogeneous multiple integral
commutes with taking its degree. Thus the same assertion holds
for every past $L^2$ coefficient by polynomial density.
Clock-chaos projection is a contraction in the full time-probability
norm. Density of predictable elementary fields and closedness of
the predictable subspace extend this action to all of $\DD_T$.
For $E$-valued fields use finite tensors with vectors of $E$,
then their density; no scenery-chaos truncation is imposed.

For the elementary field, multiply its $k$th past chaos by the
independent fresh Brownian increment. In the Gaussian product
formula all contractions with that increment vanish, since its
Cameron--Martin direction is orthogonal to the past. Its stochastic
integral therefore has degree $k+1$. Expanding the coefficient
in orthogonal degrees proves \eqref{eq:shift} on the elementary
fields. Bounded projections, the It\^o isometry and density extend
it to every field. Sectionwise orthogonality and completeness give
$\sum\Pi_ku=u$ and $\sum\|\Pi_ku\|^2=\|u\|^2$. Consequently
\eqref{eq:A} converges strongly and has norm at most one.
Apply \eqref{eq:shift} to its finite sums and pass to the limit.
The It\^o integral has zero clock mean as an $E$-valued variable,
so its $P_0$ component vanishes. This proves both equalities
in \eqref{eq:intertwine}.
\end{proof}
The predictable action is an actual section action, not an operator
defined only by matching terminal stochastic integrals. Finite
projection sums commute with the product-$L^2$ time-section
isometry. Their strong convergence has an almost-everywhere
convergent subsequence in time, which identifies the limiting
section with the corresponding sectionwise chaos series wherever
that series is already known to converge.

Applying \eqref{eq:intertwine} to Proposition~\ref{prop:finite} gives
\begin{equation}\label{eq:smfinite}
 \Smo F_{\varepsilon,T}=\J_T\Reg_TU_\varepsilon
                                  +\Smo C_{\varepsilon,T}.
\end{equation}
This identity alone does not show that either term on the right
converges: their possible cancellation has not yet been controlled.

\par\noindent Source correspondences:
\par\noindent \code{GaussianHorizonChaosIntertwining.horizonProjection_eq},
\par\noindent \code{GaussianSourceHorizonRegularization.integral_sourceProjection},
\par\noindent \code{GaussianSourceHorizonRegularization.integral_sourceRegularization},
\par\noindent \code{GaussianPredictableSourceSlices.sourceProjection_slices}, and
\par\noindent \code{IBMRegularizedStochasticIdentity.smoothing_integral}.

\section{A directional recovery estimate}
Let $H_B$ be the clock Cameron--Martin space, and let $D_B$ be the
original closed Gaussian derivative. On finite chaos expansions,
$D_B$ takes degree $j$ to degree $j-1$ and
\[
 \|D_BP_jf\|^2=j\|P_jf\|^2.
\]
It follows by orthogonal summation that $D_BN_B^{-1}$ extends
boundedly to all $\HH$, with norm at most one:
\begin{equation}\label{eq:Rbound}
 \|D_BN_B^{-1}f\|^2
     =\sum_{j\geq1}\frac1j\|P_jf\|^2\leq\|f\|^2.
\end{equation}
Here its target is $L^2(\Omega_B;H_B\otimes E)$. Tensoring
the scalar operator with $E$ is justified first on finite
orthogonal tensors and then by completion.

\begin{lemma}[Rank-one graph recovery]\label{lem:recovery}
For a unit $v\in H_B$, suppose that $f$ is in the actual closed
derivative domain and $D_Bf=v\otimes G$. Then
\begin{equation}\label{eq:recovery}
 R_vf:=\langle v,D_BN_B^{-1}f\rangle
       =\sum_{k\geq0}\frac1{k+1}P_kG,
 \qquad \|R_vf\|\leq\|f\|.
\end{equation}
The series converges strongly in $\HH$.
\end{lemma}
\begin{proof}
The bound follows from \eqref{eq:Rbound} and contraction of a
Hilbert tensor against a unit vector. The closed derivative
chaos formula gives $D_BP_{k+1}f=P_kD_Bf=v\otimes P_kG$.
In $D_BN_B^{-1}f$ this component has coefficient $1/(k+1)$.
Contracting against $v$ proves the displayed coefficient formula.
The series is strong by orthogonality and $G\in\HH$; alternatively
its partial sums converge as the projections of the bounded
recovery operator. Thus the result uses the original derivative
graph rather than merely proposed coefficients for $G$.
\end{proof}

At time $t>0$ the actual clock coordinate is
$B_t=\sqrt t\,X(v_t)$. Write the reflected endpoint as
\[
 F_{\varepsilon,t}
   =h_\varepsilon(\sqrt t\,|X(v_t)|)-h_\varepsilon(0).
\]
It has closed clock derivative
\begin{equation}\label{eq:rankgraph}
 D_BF_{\varepsilon,t}=\sqrt t\,v_t\otimes U_{\varepsilon,t}.
\end{equation}
There is no ordinary differentiability assertion at the reflection
cusp. To check the closed graph, scalarize by a continuous $\ell:E\to\R$.
The function $x\mapsto\ell(h_\varepsilon(\sqrt t|x|)-h_\varepsilon(0))$
is globally Lipschitz and bounded. Its weak derivative is
$\sqrt t\operatorname{sgn}(x)\ell(h_\varepsilon'(\sqrt t|x|))$:
integrate separately on each side of zero; continuity makes the
two boundary terms cancel. Smooth approximation and compact
cutoffs give convergence of function and derivative in Gaussian
$L^2$, using their uniform bounds. Lifting the one-dimensional
graph to $X(v_t)$ gives \eqref{eq:rankgraph} after every $\ell$.
Finite-dimensional projections in the separable value space,
and closedness of the vector derivative, recover the full graph.

Set
\begin{equation}\label{eq:qeps}
 q_\varepsilon(t)=t^{-1/2}R_{v_t}F_{\varepsilon,t},\qquad
 q_0(t)=t^{-1/2}R_{v_t}F_{0,t}\quad(t>0).
\end{equation}
Equations \eqref{eq:rankgraph} and \eqref{eq:recovery} identify
$q_\varepsilon(t)$ with the time section of $\Reg_TU_\varepsilon$
for almost every $t\in(0,T]$. No derivative-domain membership of
$F_{0,t}$ is assumed in defining $q_0$: the operator $R_{v_t}$
is already bounded on all $\HH$.

\begin{proposition}[The integrable singular bound]\label{prop:majorant}
For every $t>0$,
\begin{equation}\label{eq:qbound}
 \|q_\varepsilon(t)\|^2\leq L^2m_1t^{-1/2},\qquad
 \|q_0(t)\|^2\leq L^2m_1t^{-1/2}.
\end{equation}
For every fixed $t>0$, $q_\varepsilon(t)\to q_0(t)$ in $\HH$
as the positive width tends to zero, independently of the profile.
\end{proposition}
\begin{proof}
By the recovery bound,
$\|q_\varepsilon(t)\|^2\leq t^{-1}\|F_{\varepsilon,t}\|^2$.
Insert \eqref{eq:Fenergy}. The limit version follows in exactly
the same way from the second endpoint bound. At fixed $t$, apply
the bounded operator $t^{-1/2}R_{v_t}$ to Lemma~\ref{lem:endpoint}.
Both its argument at the limit and the operator itself are fixed
by the original sources, so the resulting limit does not depend
on a smoothing profile.
\end{proof}
This is the crucial improvement over estimating the raw slope
by $\|\rho_\varepsilon\|_2$. It estimates the recovered slope
through the endpoint energy, before integrating in time.

\par\noindent Source correspondences:
\par\noindent \code{GaussianDirectionalRecovery.shifted_derivative_projection},
\par\noindent \code{IBMCompactClockGraph.tested_closed_graph},
\par\noindent \code{IBMCompactClockRecovery.recoveredSlope_projection},
\par\noindent \code{IBMCompactClockRecovery.recoveredSlope_uniform_energy}, and
\par\noindent \code{IBMCompactClockRecovery.recoveredSlope_continuum_tendsto}.

\section{Whole-time convergence in the original predictable space}
Pointwise section convergence is insufficient without a time
majorant and a measurability argument. We supply both.

\begin{proof}[Proof of Theorem~\ref{thm:main}]
First use one fixed admissible profile and the widths
$\varepsilon_n=1/(n+1)$. Each $q_{\varepsilon_n}$ is the
time-section representative of a genuine predictable $L^2$ field.
After zero extension outside $(0,T]$, these maps are strongly
measurable from time to $\HH$. Their pointwise $\HH$ limit on
$(0,T]$ is $q_0$ by Proposition~\ref{prop:majorant}; the common
null set needed to identify countably many section representatives
is harmless. Hence $1_{(0,T]}q_0$ is almost everywhere strongly
measurable. Its squared norm is bounded by the integrable function
\[
 b_T(t)=1_{(0,T]}(t)L^2m_1t^{-1/2},\qquad
 \int_\R b_T(t)\,dt=2L^2m_1\sqrt T.
\]
Consequently it is an actual Bochner time-$L^2$ element.

Now choose any positive widths $\delta_n\to0$ and any fixed
admissible profile. For almost every $t>0$ the sections converge
to that same $q_0(t)$. The elementary squared-norm inequality gives
\[
 \|q_{\delta_n}(t)-q_0(t)\|^2
 \leq2\|q_{\delta_n}(t)\|^2+2\|q_0(t)\|^2
 \leq4L^2m_1t^{-1/2}.
\]
Dominated convergence therefore yields
\begin{equation}\label{eq:timeconv}
 \int_0^T\|q_{\delta_n}(t)-q_0(t)\|_{\HH}^2\,dt\longrightarrow0.
\end{equation}
The zero-extension and time-section map embeds $\DD_T$ isometrically
into $L^2(\R;\HH)$. Its image is closed, since $\DD_T$ is complete.
The canonical sequence from the first paragraph thus shows that
$1_{(0,T]}q_0$ is in this image. Let $u_{0,T}$ be its unique
preimage. Equation \eqref{eq:timeconv} proves convergence to it
in the original space $\DD_T$ for every such sequence of widths.
In metric spaces the positive-width limit at zero is characterized
by all positive sequences tending to zero. If that limit failed,
one could choose successively $0<\delta_n<1/(n+1)$ whose errors
stay above a fixed positive number, contradicting
\eqref{eq:timeconv}. This proves \eqref{eq:uconv} for the whole
continuum of widths, not just the canonical sequence.

The definition \eqref{eq:qzero} is independent of the profile, so
the predictable limit is independent as well. Integrating
\eqref{eq:qbound} gives \eqref{eq:energy}. By endpoint convergence
and boundedness of $\Smo$, the left side of \eqref{eq:smfinite}
converges in $\HH$. By \eqref{eq:uconv} and the It\^o isometry
its stochastic term also converges. Thus its actual remaining
term satisfies
\[
 \Smo C_{\varepsilon,T}\longrightarrow
 c_{0,T}:=\Smo F_{0,T}-\J_Tu_{0,T}
 \quad\text{in }\HH.
\]
This proves \eqref{eq:cconv} and \eqref{eq:limitidentity}.
The residual expression defines the limit of already constructed
raw corrections; it does not replace their finite-width
definition \eqref{eq:rawc}.
\end{proof}
Every scenery sector is present in this argument. The convergence
of sections is in $\HH$, and the time norm in
\eqref{eq:timeconv} integrates its full joint squared norm. A
proof only after scenery expectation or only after finitely many
source tests would not establish this result.

\par\noindent Source correspondences:
\par\noindent \code{GaussianPredictableSourceSlices.regularization_slices_of_clock_limit},
\par\noindent \code{IBMRegularizedTimeLimit.majorant_integrable},
\par\noindent \code{IBMRegularizedTimeLimit.exists_predictable_limit},
\par\noindent \code{IBMDilatedPredictableLimit.slices_sequence_tendsto},
\par\noindent \code{IBMDilatedPredictableLimit.field_continuum_tendsto}, and
\par\noindent \code{IBMDilatedPredictableLimit.correction_continuum_tendsto}.

\section{The exact correction norm and fixed-sector recovery}
The smoothing in \eqref{eq:cconv} specifies a genuine weaker
norm on the original joint variables:
\begin{equation}\label{eq:dualnorm}
 \|f\|_-:=\|\Smo f\|_{\HH}.
\end{equation}
Its definiteness follows from injectivity. It is not equivalent
to the original norm: if $e_k$ is a unit vector in clock chaos
$k\geq1$, then $\|e_k\|_-=1/k$ while $\|e_k\|=1$.

Give the actual test image $\mathcal T=\Smo\HH$ the unique
preimage norm $\|\Smo\psi\|_{\mathcal T}=\|\psi\|_{\HH}$.
This is well-defined by injectivity and makes $\mathcal T$
a Hilbert space. Its inclusion in $\HH$ is continuous and dense.
Self-adjointness and Cauchy--Schwarz give
\begin{equation}\label{eq:dualattain}
 \sup_{\|\psi\|\leq1}|\langle f,\Smo\psi\rangle|
 =\sup_{\|\psi\|\leq1}|\langle\Smo f,\psi\rangle|
 =\|f\|_-.
\end{equation}
For a nonzero smoothed field the last supremum is attained at
$\psi=\Smo f/\|\Smo f\|$; it is zero when $f=0$.
Thus this is the operator norm of an actual functional on a
specified test space, not only a collection of individual pairings.
The completion of the original space in this norm is isometric
to $\HH$ via the dense-range isometry $f\mapsto\Smo f$.
Under this identification the limit correction is represented
by $c_{0,T}$. It need not have an unsmoothed $\HH$ representative.
The completion observation is a direct human Hilbert-space
consequence of injectivity and dense range, not an extra named
formal completion construction.

\begin{proposition}[Uniform dual convergence]\label{prop:dual}
With the limits of Theorem~\ref{thm:main},
\begin{equation}\label{eq:uniformdual}
 \sup_{\|\psi\|_{\HH}\leq1}
 |\langle C_{\varepsilon,T},\Smo\psi\rangle
                  -\langle c_{0,T},\psi\rangle|\longrightarrow0.
\end{equation}
For each fixed $k\geq0$,
\begin{equation}\label{eq:fixedsector}
 P_kC_{\varepsilon,T}\longrightarrow w_k^{-1}P_kc_{0,T}
                 \quad\text{in the original }\HH.
\end{equation}
\end{proposition}
\begin{proof}
The expression under the supremum is
$|\langle\Smo C_{\varepsilon,T}-c_{0,T},\psi\rangle|$.
Its supremum is the norm of this difference, which tends to zero
by \eqref{eq:cconv}. For the fixed sector apply $P_k$ to the
same convergence, use $P_k\Smo=w_kP_k$, and multiply by the
fixed scalar $w_k^{-1}$.
\end{proof}
The inverse factor in \eqref{eq:fixedsector} is $1$ at $k=0$
and $k$ at positive $k$. It is unbounded across sectors.
Convergence of every fixed sector therefore does not justify
summing the raw limits in $\HH$.

\par\noindent Source correspondences:
\par\noindent \code{IBMCommonClockDual.smoothing_pairing},
\par\noindent \code{IBMCommonClockDual.smoothing_denseRange},
\par\noindent \code{IBMCommonClockDual.dual_attainment},
\par\noindent \code{IBMDilatedDualClosure.correction_dual_uniform}, and
\par\noindent \code{IBMDilatedDualClosure.raw_clock_projection_tendsto}.

\section{Independent rectangular source and clock tails}
Let $Q_m$ be the $m$th scenery-chaos projection, acting on $E$.
It commutes with every clock projection. Set
\[
 R_{M,N}=\Big(\sum_{m=0}^M Q_m\Big)
          \Big(\sum_{k=0}^N P_k\Big).
\]
These are orthogonal projections on $\HH$. They commute with
$\Smo$ and converge strongly to the identity as $M,N$ tend
independently to infinity. Indeed, writing $Q_{\leq M}$ and
$P_{\leq N}$ for the two factors,
\[
 \|f-R_{M,N}f\|
 \leq\|f-Q_{\leq M}f\|+\|f-P_{\leq N}f\|,
\]
which tends to zero by completeness of the separate chaos
decompositions. The two truncation parameters need not coincide.

\begin{theorem}[Uniform full-correction tails]\label{thm:tails}
For every $\eta>0$ there exist $M_0,N_0\in\mathbb N$ and
$\varepsilon_0>0$ such that for all $M\geq M_0$, $N\geq N_0$
and $0<\varepsilon<\varepsilon_0$,
\begin{equation}\label{eq:tails}
 \|C_{\varepsilon,T}-R_{M,N}C_{\varepsilon,T}\|_-<\eta.
\end{equation}
For each fixed $(m,k)$ the original joint block satisfies
\[
 Q_mP_kC_{\varepsilon,T}\longrightarrow
                    w_k^{-1}Q_mP_kc_{0,T}\quad\text{in }\HH.
\]
\end{theorem}
\begin{proof}
Choose $M_0,N_0$ so that
$\|(I-R_{M_0,N_0})c_{0,T}\|<\eta/2$.
Larger rectangular projections have smaller errors, since their
ranges contain this range and they are orthogonal projections.
Choose $\varepsilon_0$ so that
$\|\Smo C_{\varepsilon,T}-c_{0,T}\|<\eta/2$ for every
$0<\varepsilon<\varepsilon_0$. Commutation and contractivity give
\begin{align*}
 \|C_{\varepsilon,T}-R_{M,N}C_{\varepsilon,T}\|_-
 &=\|(I-R_{M,N})\Smo C_{\varepsilon,T}\|\\
 &\leq\|\Smo C_{\varepsilon,T}-c_{0,T}\|
                    +\|(I-R_{M,N})c_{0,T}\|<\eta.
\end{align*}
The fixed-block assertion follows by applying the contraction
$Q_m$ to \eqref{eq:fixedsector}.
\end{proof}
This uniform result follows from strong convergence of the full
smoothed correction. Separate convergence for every fixed block
would not imply it. Equally, \eqref{eq:tails} is a dual-norm
tail estimate; it asserts no uniform tail bound for the raw
correction in the original norm. Both qualifiers are essential.

\par\noindent Source correspondences:
\par\noindent \code{IBMDilatedDualClosure.uniform_joint_tail},
\par\noindent \code{IBMDilatedDualClosure.correction_uniform_rectangular_tail}, and
\par\noindent \code{IBMDilatedDualClosure.raw_block_tendsto}.

\section{Why the clock vacuum is retained}
A scalar model explains the $P_0$ term without appealing to the
IBM limit. Let $\gamma$ be standard Gaussian measure and let
$h_n$ be its orthonormal probabilists' Hermite polynomials.
Then $h_n'=\sqrt n\,h_{n-1}$. Define
\[
 \mathcal Bf=\sum_{n\geq0}
       w_n\sqrt{(n+1)(n+2)}\,\langle f,h_{n+2}\rangle h_n.
\]
For $n=0$ the squared gain is $2$, while for $n\geq1$ it is
$(n+1)(n+2)/n^2\leq6$. Orthogonality proves
$\|\mathcal B\|\leq\sqrt6$; equality holds on $h_3$, whose
image is $\sqrt6h_1$. If $f$ has two successive actual closed
Gaussian derivatives, differentiating its Hermite coefficients
twice gives
\[
 \mathcal Bf=\Smo f''.
\]
This equality concerns the actual closed derivative graph.
The bounded left side extends beyond that twice differentiable
domain, without asserting that the ordinary second derivative
has become an $L^2$ function there.

In particular $|x|$ has first closed derivative
$\operatorname{sgn}(x)$, and
\[
 \langle\mathcal B|x|,h_0\rangle=\E|G|>0.
\]
To check the last equality directly, $h_2=(x^2-1)/\sqrt2$,
so the constant coefficient is
$\E[|G|(G^2-1)]$. Integration on the positive half-line gives
$\E|G|^3=2\E|G|$, proving the claim. The ordinary second
derivative away from zero is zero, but its use alone would
erase this nonzero boundary channel. With $N_B^{-1}$ alone
the clock-vacuum component of the correction is also erased.
Adding $P_0$ makes the common correction norm definite and
preserves this channel.

The scalar calculation is an illustration of the shared
regularization, not an identification of the entire IBM
correction with $\mathcal B|x|$. The joint formula also has
the nonlinear scenery coefficients and all clock/source sectors.
Nor does the norm $\sqrt6$ replace the different first-order
recovery bound used in Proposition~\ref{prop:majorant}.

\par\noindent Source correspondences:
\par\noindent \code{GaussianRegularizedSecondDerivative.regularizedSecond_actual_graph},
\par\noindent \code{GaussianRegularizedSecondDerivative.regularizedSecond_norm}, and
\par\noindent \code{GaussianReflectedSecondBoundary.absolute_boundary_coefficient}.

\section{Classical context and the remaining mathematical boundaries}
Gaussian number operators, chaos projections, Hilbert-valued It\^o
integration and reflected local-time formulas are classical.
Maas and van Neerven~\cite{maas} prove vector-valued Clark--Ocone
representations in UMD spaces. Their Theorems 6.6 and 6.7 include
the differentiated and boundedly extended representations;
the latter's complete short proof uses density and the integral
isomorphism onto its range. We directly compared those statements
and their proofs, and their closed derivative/divergence setup.
This is relevant prior art for value-space stochastic integration
and extension by completion. Our source space $E=L^2(\Omega_W)$
is a Hilbert space, so its vector-valued nature is not a new
stochastic-integration principle. We did not review every proof
or hypothesis of their general UMD integration theorem.

The specific result assembled here concerns the actual reflected
scenery family, its mollifier-independent recovered integrand,
the integrable origin bound and the common correction topology.
The bounded literature search does not certify priority for that
assembly. This is a formalization companion with complete proofs
of its displayed convergence mechanism and exact formal source
correspondences, not a claim to a newly discovered Gaussian calculus.

Three distinctions delimit the result. First, $\J_T$ integrates
against the original Brownian clock. The entire scenery remains
the value-space parameter, so the integrand may use information
unavailable from the observed past of $Z$. No theorem identifies
$u_{0,T}$ as an observed-filtration IBM integrand.

Second, $\Smo$ is fixed and injective but has no bounded inverse
on its range with the inherited $\HH$ norm. The high-chaos
vectors following \eqref{eq:dualnorm} prove this directly.
One cannot apply an uncontrolled inverse to
\eqref{eq:limitidentity} to obtain an unsmoothed $L^2$ It\^o
identity. The theorem does not prove that the raw slopes or
raw corrections converge, and does not assert they diverge.
The two terms in \eqref{eq:rawc} are only controlled jointly
after the common smoothing in the limit theorem.

Third, the convergence is strong in the entire original
predictable norm for the regularized integrand, strong in the
common dual norm for the full correction, and strong in ordinary
joint $L^2$ for each fixed clock sector. These different claims
cannot be substituted for one another. Neither independent
rectangular dual tails nor full source retention supplies
unsmoothed raw $L^2$ tails.

The other IBM companions~\cite{companions} treat spatial grid
contractions, faithful random-kernel test spaces, graph-domain
obstructions, actual time averages and static intrinsic adjoints.
Those statements retain their own topologies. In particular the
present regularization does not remove the positive-time native
clock derivative obstruction for IBM point observations. A full
physical-time calculus in the observed filtration would still
require a specified closed observed integral, admissible causal
coefficients and a reconstruction theorem in its own topology.
Those obligations remain open here.

\begin{thebibliography}{9}
\bibitem{maas} J. Maas and J. van Neerven,
\emph{A Clark--Ocone formula in UMD Banach spaces},
Electronic Communications in Probability \textbf{13} (2008), 151--164.
Primary author manuscript, arXiv:0709.2021v2, Sections 3--4
and Theorems 6.6--6.7 (statements and the displayed proofs).
\url{https://arxiv.org/pdf/0709.2021}.
\bibitem{source} Project Stardust, \emph{OperatorDerivativeRewrite},
frozen verified source commit given in Section 1. The accompanying
\texttt{theorem-map.json} records exact headers, hashes and line links.
\bibitem{companions} Project Stardust, reviewed editorial companions:
\emph{Variation and spatial Simpson integration for reflected-clock iterated Brownian motion}
(\texttt{../ibm-spatial-calculus/main.pdf});
\emph{Gaussian clock regularity and domain obstructions for reflected-clock Brownian motion}
(\texttt{../ibm-clock-domain/main.pdf});
\emph{Brownian clock mixtures}
(\texttt{../ibm-one-time-laws/main.pdf}); and the constructed
Tanaka/holding-measure account in
\texttt{../sticky-clock-recovery/main.pdf}.
These are independent publication drafts from the same frozen baseline,
not external priority evidence.
\end{thebibliography}
\end{document}
